% Options for packages loaded elsewhere
\PassOptionsToPackage{unicode}{hyperref}
\PassOptionsToPackage{hyphens}{url}
\documentclass[
  ignorenonframetext,
]{beamer}
\newif\ifbibliography
\usepackage{pgfpages}
\setbeamertemplate{caption}[numbered]
\setbeamertemplate{caption label separator}{: }
\setbeamercolor{caption name}{fg=normal text.fg}
\beamertemplatenavigationsymbolsempty
% remove section numbering
\setbeamertemplate{part page}{
  \centering
  \begin{beamercolorbox}[sep=16pt,center]{part title}
    \usebeamerfont{part title}\insertpart\par
  \end{beamercolorbox}
}
\setbeamertemplate{section page}{
  \centering
  \begin{beamercolorbox}[sep=12pt,center]{section title}
    \usebeamerfont{section title}\insertsection\par
  \end{beamercolorbox}
}
\setbeamertemplate{subsection page}{
  \centering
  \begin{beamercolorbox}[sep=8pt,center]{subsection title}
    \usebeamerfont{subsection title}\insertsubsection\par
  \end{beamercolorbox}
}
% Prevent slide breaks in the middle of a paragraph
\widowpenalties 1 10000
\raggedbottom
\AtBeginPart{
  \frame{\partpage}
}
\AtBeginSection{
  \ifbibliography
  \else
    \frame{\sectionpage}
  \fi
}
\AtBeginSubsection{
  \frame{\subsectionpage}
}
\usepackage{iftex}
\ifPDFTeX
  \usepackage[T1]{fontenc}
  \usepackage[utf8]{inputenc}
  \usepackage{textcomp} % provide euro and other symbols
\else % if luatex or xetex
  \usepackage{unicode-math} % this also loads fontspec
  \defaultfontfeatures{Scale=MatchLowercase}
  \defaultfontfeatures[\rmfamily]{Ligatures=TeX,Scale=1}
\fi
\usepackage{lmodern}
\ifPDFTeX\else
  % xetex/luatex font selection
\fi
% Use upquote if available, for straight quotes in verbatim environments
\IfFileExists{upquote.sty}{\usepackage{upquote}}{}
\IfFileExists{microtype.sty}{% use microtype if available
  \usepackage[]{microtype}
  \UseMicrotypeSet[protrusion]{basicmath} % disable protrusion for tt fonts
}{}
\makeatletter
\@ifundefined{KOMAClassName}{% if non-KOMA class
  \IfFileExists{parskip.sty}{%
    \usepackage{parskip}
  }{% else
    \setlength{\parindent}{0pt}
    \setlength{\parskip}{6pt plus 2pt minus 1pt}}
}{% if KOMA class
  \KOMAoptions{parskip=half}}
\makeatother
\usepackage{longtable,booktabs,array}
\newcounter{none} % for unnumbered tables
\usepackage{calc} % for calculating minipage widths
\usepackage{caption}
% Make caption package work with longtable
\makeatletter
\def\fnum@table{\tablename~\thetable}
\makeatother
\setlength{\emergencystretch}{3em} % prevent overfull lines
\providecommand{\tightlist}{%
  \setlength{\itemsep}{0pt}\setlength{\parskip}{0pt}}
\usepackage{bookmark}
\IfFileExists{xurl.sty}{\usepackage{xurl}}{} % add URL line breaks if available
\urlstyle{same}
\hypersetup{
  hidelinks,
  pdfcreator={LaTeX via pandoc}}

\author{\texorpdfstring{}{}}
\date{}

\begin{document}

\begin{frame}
\newpage
\end{frame}

\begin{frame}{Attack Corpus: Methodology and Ethical Considerations}
\protect\phantomsection\label{sec:attack-methodology}
This section documents the attack generation methodology, effectiveness
analysis, ethical considerations, and data availability.

\begin{block}{Attack Generation Methodology}
\protect\phantomsection\label{sec:generation-methodology}
\begin{block}{Synthetic Attack Generation}
\protect\phantomsection\label{sec:synthetic-generation}
\textbf{Process}:

\begin{enumerate}
\item **Template Creation**: Define attack structure templates for each category
\item **Parameter Variation**: Systematically vary attack parameters
\item **Constraint Satisfaction**: Ensure attacks satisfy category definitions
\item **Deduplication**: Remove semantically equivalent attacks
\item **Validation**: Human review of generated attacks
\end{enumerate}

\textbf{Table: Generation method statistics.} \{\#tab:generation-stats\}

{\def\LTcaptype{none} % do not increment counter
\begin{longtable}[]{@{}llll@{}}
\toprule\noalign{}
Method & Count & QA Pass Rate\(^*\) & Mean Sophistication \\
\midrule\noalign{}
\endhead
Template instantiation & 420 & 91\% & 0.4 \\
Parameter variation & 250 & 88\% & 0.5 \\
Red team manual crafting & 150 & 95\% & 0.8 \\
LLM-assisted mutation & 80 & 75\% & 0.7 \\
Adversarial optimization & 50 & 82\% & 0.9 \\
\bottomrule\noalign{}
\end{longtable}
}

\(^*\)\textit{QA pass rate denotes the proportion of candidate attacks that passed quality assurance validation (measurability, reproducibility, category alignment), not attack efficacy against defended systems. Total candidates generated exceeded 1,200; 950 passed QA and were retained.}
\end{block}

\begin{block}{Red Team Exercise Protocol}
\protect\phantomsection\label{sec:red-team}
\textbf{Participants}: 8 security researchers (2--10 years experience)

\textbf{Duration}: 4 weeks

\textbf{Methodology}:

\begin{enumerate}
\item **Week 1**: Familiarization with target architectures
\item **Week 2**: Independent attack development
\item **Week 3**: Cross-team attack validation
\item **Week 4**: Documentation and categorization
\end{enumerate}
\end{block}

\begin{block}{Quality Assurance}
\protect\phantomsection\label{sec:qa}
\textbf{Table: Attack validation criteria.}
\{\#tab:validation-criteria\}

{\def\LTcaptype{none} % do not increment counter
\begin{longtable}[]{@{}ll@{}}
\toprule\noalign{}
Criterion & Description \\
\midrule\noalign{}
\endhead
Measurability & Success/failure unambiguously determinable \\
Reproducibility & Attack produces consistent results \\
Category alignment & Attack matches labeled category \\
Non-trivial & Attack not detected by simple heuristics \\
\bottomrule\noalign{}
\end{longtable}
}

\textbf{Validation Process}:

\begin{enumerate}
\item Two independent reviewers per attack
\item Disagreements resolved by third reviewer
\item Inter-rater reliability: Cohen's $\kappa = 0.84$
\end{enumerate}
\end{block}
\end{block}

\begin{block}{Attack Effectiveness Analysis}
\protect\phantomsection\label{sec:effectiveness-analysis}
\begin{block}{Success Rate by Defense Configuration}
\protect\phantomsection\label{sec:success-by-defense}
\textbf{Table: Attack success rate by defense configuration.}
\{\#tab:success-by-defense\}

{\def\LTcaptype{none} % do not increment counter
\begin{longtable}[]{@{}lllll@{}}
\toprule\noalign{}
Defense & Prompt Inj. & Trust Expl. & Belief Manip. & Coord. \\
\midrule\noalign{}
\endhead
Firewall only & 15\% & 38\% & 29\% & 42\% \\
Sandbox only & 35\% & 25\% & 31\% & 55\% \\
Tripwires only & 22\% & 18\% & 8\% & 48\% \\
Full CIF & 4\% & 9\% & 7\% & 11\% \\
\bottomrule\noalign{}
\end{longtable}
}
\end{block}

\begin{block}{Attack Sophistication Correlation}
\protect\phantomsection\label{sec:sophistication-corr}
\begin{equation}
\label{eq:sophistication-correlation}
\rho_{sophistication, success} = 0.67 \quad (p < 0.001, \; n = 950)
\end{equation}

We report Spearman's rank correlation (\(\rho\)) rather than Pearson's
\(r\) because sophistication levels (Low, Medium, High, Expert) are
ordinal categories. More sophisticated attacks have higher baseline
success but show similar detection rates under CIF, suggesting defense
robustness.
\end{block}

\begin{block}{Temporal Analysis}
\protect\phantomsection\label{sec:temporal-analysis}
\textbf{Table: Detection rate by attack age.} \{\#tab:attack-age\}

{\def\LTcaptype{none} % do not increment counter
\begin{longtable}[]{@{}lll@{}}
\toprule\noalign{}
Attack Age & Detection Rate & \(n\) \\
\midrule\noalign{}
\endhead
\(<\) 6 months & 91\% & 285 \\
6--12 months & 94\% & 380 \\
\(>\) 12 months & 96\% & 285 \\
\bottomrule\noalign{}
\end{longtable}
}

Older attacks are detected at higher rates due to pattern database
inclusion. The 5-point gap between newest and oldest cohorts quantifies
the advantage that known-signature detection provides and underscores
the importance of continuous corpus expansion to maintain efficacy
against novel techniques.
\end{block}
\end{block}

\begin{block}{Ethical Considerations}
\protect\phantomsection\label{sec:ethical-considerations}
\begin{block}{Responsible Disclosure}
\protect\phantomsection\label{sec:responsible-disclosure}
All novel attack vectors discovered during this research were:

\begin{enumerate}
\item **Reported**: Communicated to affected framework maintainers
\item **Embargoed**: 90-day disclosure window before publication
\item **Mitigated**: Defenses provided alongside vulnerability reports
\end{enumerate}

\textbf{Table: Disclosure timeline.} \{\#tab:disclosure-timeline\}

{\def\LTcaptype{none} % do not increment counter
\begin{longtable}[]{@{}
  >{\raggedright\arraybackslash}p{(\linewidth - 6\tabcolsep) * \real{0.2500}}
  >{\raggedright\arraybackslash}p{(\linewidth - 6\tabcolsep) * \real{0.2500}}
  >{\raggedright\arraybackslash}p{(\linewidth - 6\tabcolsep) * \real{0.2500}}
  >{\raggedright\arraybackslash}p{(\linewidth - 6\tabcolsep) * \real{0.2500}}@{}}
\toprule\noalign{}
\begin{minipage}[b]{\linewidth}\raggedright
Framework
\end{minipage} & \begin{minipage}[b]{\linewidth}\raggedright
Date Reported
\end{minipage} & \begin{minipage}[b]{\linewidth}\raggedright
Status
\end{minipage} & \begin{minipage}[b]{\linewidth}\raggedright
Resolution
\end{minipage} \\
\midrule\noalign{}
\endhead
Framework A & 2025-06-15 & Acknowledged & Patched (v2.1.3) \\
Framework B & 2025-06-22 & Acknowledged & In progress \\
Framework C & 2025-07-01 & No response & Public disclosure (90-day
window elapsed) \\
Framework D & 2025-07-10 & Acknowledged & Mitigated via configuration
change \\
\bottomrule\noalign{}
\end{longtable}
}

Framework names are anonymized per coordinated disclosure agreements.
Specific vulnerability details are available to affected maintainers and
will be published after all embargo periods expire.
\end{block}

\begin{block}{Dual-Use Considerations}
\protect\phantomsection\label{sec:dual-use}
\textbf{Risk Assessment}: The attack corpus represents a dual-use
resource that could enable both defensive research and malicious
exploitation. We address this through:

\begin{enumerate}
\item **Sanitization**: All published examples are non-functional
\item **Partial Disclosure**: Full corpus available only to verified researchers
\item **Access Controls**: Request-based access with institutional verification
\item **Usage Tracking**: Audit log of corpus access
\end{enumerate}

\textbf{Table: Access control hierarchy.} \{\#tab:access-hierarchy\}

{\def\LTcaptype{none} % do not increment counter
\begin{longtable}[]{@{}lll@{}}
\toprule\noalign{}
Access Level & Scope & Requirement \\
\midrule\noalign{}
\endhead
Researcher & Template structures & Institutional affiliation \\
Full access & Complete corpus & IRB approval + NDA \\
\bottomrule\noalign{}
\end{longtable}
}
\end{block}

\begin{block}{Defense Framework Dual-Use Considerations}
\protect\phantomsection\label{sec:defense-dual-use}
While the attack corpus dual-use considerations are addressed above, the
defense framework itself presents distinct dual-use risks that warrant
separate analysis.

\textbf{Detection Algorithm Inversion Risk.} The detection algorithms
documented in \cref{sec:detection-algorithms} could potentially be
analyzed to design evasive attacks that remain below detection
thresholds. An adversary with access to the full algorithm
specifications could craft attacks that exploit known blind spots or
systematically probe the feature space to identify classification
boundaries. This risk is inherent to any published detection
methodology.

\textbf{Trust Calculus Parameter Exposure.} The trust decay parameter
(\(\delta\)), delegation depth limits, and threshold configurations
disclosed in this paper could enable adversaries to game the trust
system if they know the exact values deployed in a target system.
Attackers could craft delegation chains that remain just above trust
thresholds or time their attacks to coincide with trust recovery
periods.

\textbf{Observed Mitigation Approaches.} Several approaches address
these dual-use risks:

\begin{enumerate}
\item **API Abstraction**: Deploying CIF through an abstraction layer that hides internal parameters. Detection decisions exposed as binary outcomes (allowed/blocked) without revealing confidence scores or feature contributions.
\item **Parameter Randomization**: Introducing slight randomization in threshold values and decay parameters across instances, reducing the exploitability of published defaults.
\item **Adversarial Probing Detection**: Monitoring for patterns indicative of boundary probing (repeated near-threshold submissions, systematic parameter variation).
\end{enumerate}

The defense composition algebra (established in Part 1) remains valid
regardless of specific parameter choices, ensuring that the theoretical
guarantees hold even when operational parameters differ from published
defaults. Specific deployment configurations are detailed in Part 3.
\end{block}

\begin{block}{Human Subjects}
\protect\phantomsection\label{sec:human-subjects}
This research did not involve human subjects experimentation. All
attacks were tested against:

\begin{itemize}
\item Synthetic agent configurations
\item Sandboxed environments
\item No production systems with real users
\end{itemize}
\end{block}

\begin{block}{Research Ethics Approval}
\protect\phantomsection\label{sec:ethics-approval}
This research was reviewed and determined to be exempt from IRB
oversight as it did not involve human subjects. The board determined
that:

\begin{enumerate}
\item No human subjects were involved
\item Dual-use risks were adequately mitigated
\item Responsible disclosure practices were followed
\end{enumerate}
\end{block}
\end{block}

\begin{block}{Data Availability}
\protect\phantomsection\label{sec:data-availability}
\begin{block}{Public Resources}
\protect\phantomsection\label{sec:public-resources}
\begin{itemize}
\item Sanitized attack examples: This supplementary material
\item Detection patterns: Available in paper repository
\item Defense implementations: Available at DOI: 10.5281/zenodo.18364128
\end{itemize}
\end{block}

\begin{block}{Restricted Resources}
\protect\phantomsection\label{sec:restricted-resources}
\begin{itemize}
\item Full attack corpus: Available upon request
\item Red team exercise data: Institution members only
\item Unpublished vulnerabilities: Covered by disclosure agreements
\end{itemize}
\end{block}

\begin{block}{Access Request Process}
\protect\phantomsection\label{sec:access-request}
Researchers wishing to access the full attack corpus must:

\begin{enumerate}
\item Submit institutional affiliation verification
\item Provide IRB approval or exemption letter
\item Sign data use agreement
\item Agree to responsible use terms
\end{enumerate}
\end{block}
\end{block}

\begin{block}{References}
\protect\phantomsection\label{sec:corpus-references}
The attack corpus draws from JailbreakBench
\cite{chao2024jailbreakbench}, PromptInject \cite{liu2023prompt}, and
TensorTrust \cite{toyer2024tensortrust}.
\end{block}
\end{frame}

\end{document}
